% Generated draft: requires mathematical and source review.
\documentclass[11pt]{article}
\usepackage[margin=25mm]{geometry}
\usepackage{fontspec}
\setmainfont{Latin Modern Roman}
\usepackage{amsmath,amssymb,mathtools,unicode-math}
\setmathfont{Latin Modern Math}
\usepackage{enumitem,needspace,xurl,hyperref}
\setlength{\parindent}{0pt}
\setlength{\parskip}{5pt}
\setlength{\emergencystretch}{3em}
\newcommand{\N}{\mathbb{N}}
\newcommand{\Z}{\mathbb{Z}}
\newcommand{\Q}{\mathbb{Q}}
\newcommand{\R}{\mathbb{R}}
\newcommand{\C}{\mathbb{C}}
\newcommand{\F}{\mathbb{F}}
\newcommand{\A}{\mathbb{A}}
\newcommand{\PP}{\mathbb{P}}
\newcommand{\E}{\mathbb{E}}
\newcommand{\eps}{\varepsilon}
\newcommand{\dd}{\,\mathrm{d}}
\begin{document}
\section*{UnsolvedMath ID: 30006497}
\textbf{Source problem code:} catalogue entry
\subsection{Definitions and mathematical statement}
Under the conventional interpretation of the braid category, let \(\mathcal B\) be the strict monoidal category whose objects are integers \(n\geq 0\), with \(\operatorname{Hom}_{\mathcal B}(m,n)=\varnothing\) for \(m\neq n\) and \(\operatorname{End}_{\mathcal B}(n)=B_n\), the group of isotopy classes of braids on \(n\) strands. Composition is stacking, the tensor product is juxtaposition, and the tensor unit is \(0\); \(B_0\) and \(B_1\) are trivial. The supplied record does not explicitly specify these conventions.

Let \(\mathrm{Mat}\) denote the monoidal matrix category intended by the source, and fix a monoidal subcategory \(\mathcal Y\subseteq\mathrm{Mat}\). The coefficient ring or field, the objects and allowed matrices, and the tensor product and strictness conventions of \(\mathrm{Mat}\) are not specified in the supplied record. Consequently, they remain parameters requiring clarification; no particular matrix category is selected here. The original passage assumes only a subcategory \(Y\), whereas the supplied corrected statement requires a monoidal subcategory. This draft follows the corrected statement, subject to that qualification.

Assuming compatible strict monoidal structures have been specified, a strict monoidal functor \(F:\mathcal B\to\mathcal Y\) is a functor satisfying
\[
F(0)=\mathbb 1_{\mathcal Y},\qquad
F(m+n)=F(m)\otimes F(n),\qquad
F(f\otimes g)=F(f)\otimes F(g)
\]
as equalities, in addition to preserving identities and composition. Write \(\operatorname{MonFun}(\mathcal B,\mathcal Y)\) for the source's collection of such functors; the supplied excerpt does not define the morphisms of this functor category.

\textbf{Classification problem.} For each fixed admissible \(\mathcal Y\), classify the strict monoidal functors \(F:\mathcal B\to\mathcal Y\), up to the ``appropriate equivalences'' intended by the source. The equivalence relation is not specified. Thus a precise classification target cannot be fixed from this record alone: the ambient matrix category, its monoidal conventions, and the equivalence relation must first be supplied. In particular, equivalence by monoidal natural isomorphism is not silently imposed.

\textbf{Supplied contextual qualifications.} The surrounding passage labels as Theorem 2 the assertion that \(\operatorname{MonFun}(\mathcal B,\mathrm{Mat})\) is strict monoidal under a ``lashing product,'' and mentions natural notions of subobjects and quotient objects. The product and these notions are not defined in the supplied passage, which ends before the promised subsequent properties. These are retained as source context, not additional conclusions established by this draft.
\subsection{Short English statement}
For a fixed monoidal subcategory \(\mathcal Y\) of the intended matrix category, classify strict monoidal functors from the braid category to \(\mathcal Y\). The matrix conventions and the intended equivalence relation are missing, so the classification target remains conditional. The supplied record labels the problem open; this statement-only draft requires human mathematical review.
\subsection{Sources}
\begin{itemize}
\item Supplied local UnsolvedMath registry record, permanent ID \(30006497\), title ``Classification of Monoidal Functors from the Braid Category,'' source problem code \texttt{catalogue entry}. The record identifies the extracted question as Problem 1 and gives proposed year 2025; no proposer is named.
\item Supplied source citation: \emph{Mini-Workshop: The Yang--Baxter Equation and Representations of Braid Groups (2026)}, Oberwolfach Reports, pp.~2633--2672, DOI \texttt{10.4171/owr/2025/49}. \url{https://doi.org/10.4171/owr/2025/49}. The supplied literature-source label is ``Oberwolfach Report 49/2025.'' The citation's 2026 designation and its 2025 report designation are retained without resolving their relationship.
\item Additional supplied literature source: ``Classification of spin-chain braid representations.'' \url{https://arxiv.org/abs/2112.04533}. The record provides no authors or detailed statement of its relevance or scope.
\item Literature assessment dated exactly \texttt{2026-08-22}: the supplied curation labels the problem open and describes classification for an arbitrary fixed matrix subcategory as an organizational problem not solved in general. This is supplied dated curation, not independently verified literature assessment. The background explicitly requests independent verification of time-sensitive status and cited claims.
\item The supplied registry labels the statement \texttt{corrected\_verified} and reports that the official Oberwolfach report was checked, the question isolated, and its notation repaired. These are supplied verification claims, not independent verification performed for this draft.
\end{itemize}
\end{document}
